\documentclass[11pt]{article}

\usepackage{amsmath}
\usepackage{amssymb}
\usepackage{array}
\usepackage{bm}
\usepackage{esvect}
\usepackage{geometry}
\usepackage{makecell}
\usepackage{mathtools}
\usepackage{siunitx}

\usepackage[T1]{fontenc}
\usepackage[utf8]{inputenc}

\geometry{margin=1in}

\DeclareRobustCommand{\myvec}[1]{\vv{\mathbf{#1}}}
\newcommand{\myunitvec}[1]{\bm{\hat{\mathbf{#1}}}}
\newcommand{\myvecsp}[1]{\bm{\mathcal{#1}}}
\newcommand{\mysubvecsp}[1]{\bm{{#1}}}
\newcommand{\myset}[1]{\mathcal{#1}}
\newcommand{\mymat}[1]{#1}
\newcommand{\dotprod}[2]{#1 \cdot #2}
\newcommand{\innerprod}[2]{\left\langle #1, #2\right\rangle}

\raggedright
\begin{document}
    \subsubsection*{Gram-Schmidt Orthogonalization Process}
    \begin{quote}
        Let $\{\myvec{w}_1, \myvec{w}_2, \dots, \myvec{w}_n\}$ be a basis for any vector space $\mysubvecsp{V} \subseteq \mathbb{R}^n$, then\\
        the orthogonal basis $\{\myvec{v}_1, \myvec{v}_2, \dots, \myvec{v}_n\}$ for $\mysubvecsp{V}$ can be found as follows:
    
        \begin{align*}
            \myvec{v}_1 &= \myvec{w}_1 \\
            \myvec{v}_2 &= \myvec{w}_2 - \frac{\dotprod{\myvec{v}_1}{\myvec{w}_2}}{\dotprod{\myvec{v}_1}{\myvec{v}_1}}\myvec{v}_1 \\
            \myvec{v}_3 &= \myvec{w}_3 - \frac{\dotprod{\myvec{v}_1}{\myvec{w}_3}}{\dotprod{\myvec{v}_1}{\myvec{v}_1}}\myvec{v}_1 - \frac{\dotprod{\myvec{v}_2}{\myvec{w}_3}}{\dotprod{\myvec{v}_2}{\myvec{v}_2}}\myvec{v}_2 \\
            \dots \\
            \myvec{v}_n &= \myvec{w}_n - \sum_{k = 1}^{n-1} \frac{\dotprod{\myvec{v}_k}{\myvec{w}_n}}{\dotprod{\myvec{v}_k}{\myvec{v}_k}}\myvec{v}_k \\
        \end{align*}
    \end{quote}
    \subsubsection*{Diagonalization}
    \[ \mymat{D} = \mymat{P}^{-1}\mymat{A}\mymat{P} \qquad \mymat{A} = \mymat{P}\mymat{D}\mymat{P}^{-1} \]
    
\end{document}